\subsection{Invariant functions on tubes and complex tori}
\label{sec:invariant-psh-convexity}
The complex Hessian becomes an ordinary real Hessian when a potential is independent of imaginary translations. Local integrability then excludes poles on the real base. This gives the analytic starting point for logarithmic coordinates without a separate theorem about circular means.

\begin{proposition}\label{prop:invariant-psh-convexity}
Let $V\subset\mathbb R^n$ be nonempty, open and convex, $n\geq1$, and let $U\in\PSH(V+\mathrm i\mathbb R^n)$ be invariant under imaginary translations. There is a finite convex function $f:V\to\mathbb R$ such that $U(t+\mathrm is)=f(t)$. Conversely every finite convex function on $V$ defines such a plurisubharmonic function.
\end{proposition}
\begin{proof}
Invariance defines $f(t)=U(t)$. By \cref{thm:psh-distribution-representative}, $U$ is locally integrable and finite almost everywhere. Fubini on compact real boxes times compact imaginary boxes shows that $f$ is finite almost everywhere on $V$.

On a relatively compact convex subset of $V$, radial convolution gives smooth decreasing psh approximations $U_\varepsilon$ with pointwise limit $U$. They remain independent of $s$, so write $U_\varepsilon(t+\mathrm is)=f_\varepsilon(t)$. The complex Hessian has entries
\[
 \frac{\partial^2U_\varepsilon}{\partial w_j\partial\bar w_k}
             =\frac14\frac{\partial^2f_\varepsilon}{\partial t_j\partial t_k}.
\]
Positivity makes $f_\varepsilon$ convex. Suppose $f(t_0)=-\infty$ and choose $B(t_0,2r)\Subset V$. Local upper boundedness of $U$ gives a common upper bound $M$ for the smoothings on this ball when $\varepsilon$ is small. For $t\in B(t_0,r)$, convexity gives
\[
 f_\varepsilon(t)
   \leq\tfrac12f_\varepsilon(t_0)
                     +\tfrac12f_\varepsilon(2t-t_0)
   \leq\tfrac12f_\varepsilon(t_0)+\tfrac12M.
\]
Letting $\varepsilon\downarrow0$ would make $f=-\infty$ throughout that ball, contradicting almost-everywhere finiteness. Thus $f$ is finite everywhere. Passing the convex inequalities to the pointwise limit on each compact segment proves convexity on $V$.

Conversely a finite convex $f$ is continuous and has a supporting slope at each point by \cref{prop:convex-local-regularity-and-subgradients}. At $t_0\in V$ take such a slope $p$. On every affine complex-line circle centered at $w_0=t_0+\mathrm is_0$ and contained in the tube,
\[
 f\bigl(\Real(w_0+r\mathrm e^{\mathrm i\theta}a)\bigr)
 \geq f(t_0)+\left\langle p,
                       \Real(r\mathrm e^{\mathrm i\theta}a)\right\rangle.
\]
The last term has circle mean zero. Hence $f(\Real w)$ satisfies the line submean inequality and is continuous and finite. The subharmonic characterization on each line, followed by \cref{def:plurisubharmonic}, proves the converse.
\end{proof}

\begin{corollary}\label{cor:toric-psh-finiteness}
For the same $V$, put
\[
 \Omega_V=\{z\in(\mathbb C^*)^n:
                   (\log|z_1|^2,\ldots,\log|z_n|^2)\in V\}.
\]
A rotation-invariant $\varphi\in\PSH(\Omega_V)$ has the form
\[
 \varphi(z)=f(\log|z_1|^2,\ldots,\log|z_n|^2)
\]
for a finite convex $f$ on $V$, and every such $f$ gives a rotation-invariant psh function. In particular a rotation-invariant psh function on $(\Delta^*)^n$ or $(\mathbb C^*)^n$ is finite everywhere.
\end{corollary}
\begin{proof}
The holomorphic covering $w\mapsto(\exp(w_1/2),\ldots,\exp(w_n/2))$ maps the tube onto $\Omega_V$, and its logarithmic radii are $\Real w$. Its pullback of $\varphi$ is nontrivial and psh by \cref{prop:psh-holomorphic-pullback}. Rotation invariance makes that pullback independent of imaginary translations. The proposition gives $f$ and its finiteness. Conversely the proposition gives $f(\Real w)$ psh; local holomorphic logarithms transfer this property to $\Omega_V$. The two special domains correspond to $V=(-\infty,0)^n$ and $V=\mathbb R^n$. In dimension zero the analogous assertion is the finite value at a point.
\end{proof}
\index{plurisubharmonic function!rotation invariance}
\index{convex function!invariant psh correspondence}
